\documentclass[11pt,a4paper]{article}
\usepackage[T1]{fontenc}
\usepackage[utf8]{inputenc}
\usepackage[english]{babel}
\usepackage{lmodern,microtype,amsmath,amssymb,mathtools,graphicx}
\usepackage[a4paper,left=22mm,right=22mm,top=22mm,bottom=22mm]{geometry}
\usepackage{booktabs,array,tabularx,fancyhdr,titlesec,xcolor}
\usepackage{hyperref}
\definecolor{navy}{HTML}{1E4F6D}
\definecolor{muted}{HTML}{5D7182}
\definecolor{pale}{HTML}{EAF0F5}
\hypersetup{colorlinks=true,linkcolor=navy,urlcolor=navy,
  pdftitle={Pratt Polynomials as an Argument of the Zeta Function},
  pdfsubject={Zeros, height mapping, and equivalent formulations of the Riemann Hypothesis}}
\setlength{\parindent}{0pt}
\setlength{\parskip}{5.5pt}
\setlength{\headheight}{15pt}
\setlength{\emergencystretch}{1.5em}
\renewcommand{\arraystretch}{1.16}
\pagestyle{fancy}
\fancyhf{}
\fancyhead[L]{\small\color{muted}Pratt Polynomials \textperiodcentered\ Composition with the Zeta Function}
\fancyfoot[R]{\small\color{muted}\thepage}
\renewcommand{\headrulewidth}{.3pt}
\titleformat{\section}{\Large\bfseries\color{navy}}{\thesection.}{.55em}{}
\titlespacing*{\section}{0pt}{10pt}{6pt}
\newcommand{\Zcal}{\mathcal Z}
\newcommand{\Xcal}{\mathcal X}
\begin{document}
\thispagestyle{fancy}

{\fontsize{27}{31}\selectfont\bfseries\color{navy}
Pratt Polynomials as an Argument of the Zeta Function\par}
\vspace{8pt}
{\large\bfseries Zeros, height mapping, and two equivalent forms of the Riemann Hypothesis\par}
\vspace{7pt}
{\large 5 October 2026 \textperiodcentered\ Supplement to the study of real and complex indices\par}

The proposed substitution is mathematically well defined and yields exact statements about the preimages of zeta zeros. In the critical strip, for positive real $x$, real parts are preserved. This gives an equivalent formulation of the Riemann Hypothesis, but not a proof. For a fixed complex index, concrete polynomial equations for the zeros in $x$ also arise. The derivations below proceed directly from the given family; no claim of novelty in the literature is made.

\section{Definition and Meaning of the Substitution}
We use $f_1=1$, $f_2=x$, $f_p=1+f_{p-1}$ for odd primes and $f_{ab}=f_af_b$ for positive integer indices. In addition, $f_0=0$, $f_{-n}=-f_n$, and the linear interpolation are used:
\[
 f_r(x)=(1-t)f_n(x)+t f_{n+1}(x),\qquad n=\lfloor r\rfloor,\quad t=r-n.
\]
For $s=a+ib$, the complex definition requested by the user is
\[
 f_s(x)=f_a(x)+i f_b(x).
\]
We have $f_r=r$ for $|r|\le1$, $f_{-r}=-f_r$, and $f_s(2)=s$. The new function is therefore
\[
 \Zcal(s,x)=\zeta(f_s(x)),\qquad \Zcal(s,2)=\zeta(s).
\]
Here the zeta function is understood via its meromorphic continuation. Its Dirichlet series and Euler product may be used only in the region $\Re f_s(x)>1$ [1]. In particular, no zero argument in the critical strip can be obtained from the Dirichlet series, which does not converge there.

The two variables must be kept separate: for fixed index $s$, $f_s$ is a polynomial in $x$; for fixed $x$, the map $s\mapsto f_s(x)$ is piecewise affine in the real and imaginary parts.

\section{The Real Parts Are Preserved in the Critical Strip}
\textbf{Theorem 1. Exact correspondence of the nontrivial zeros.}
For a fixed real $x$, set $h_x(t)=f_t(x)$. If $s=\sigma+it$ with $0<\sigma<1$, then $f_\sigma(x)=\sigma$ and $h_x(t)$ is real. Hence
\[
 \Zcal(\sigma+it,x)=\zeta\!\left(\sigma+i h_x(t)\right).
\]
If $\rho=\beta+i\gamma$ is a nontrivial zeta zero, then its preimages in the critical strip are exactly the points
\[
 s=\beta+it,\qquad h_x(t)=\gamma.
\]

\textbf{Proof.} The real part of the inserted zeta argument is exactly $\sigma$. Thus the equation $f_s(x)=\rho$ is equivalent to $\sigma=\beta$ and $h_x(t)=\gamma$. In the critical strip there are only nontrivial zeta zeros [2]. This makes the correspondence complete.

\textbf{Theorem 2. Every zeta zero has such a preimage for $x>0$.}
For every fixed real $x>0$, the map $h_x:\mathbb R\to\mathbb R$ is surjective.

\textbf{Proof.} The interpolation is continuous, odd, and satisfies $h_x(0)=0$. Since $f_3=x+1$ and the integer family is multiplicative,
\[
 h_x(3^k)=(x+1)^k\longrightarrow +\infty.
\]
The intermediate value theorem on $[0,3^k]$ gives every positive level, and oddness gives every negative level. Monotonicity is not needed.

\textbf{Corollary 3. An exact equivalence to the Riemann Hypothesis.}
For every arbitrarily fixed real $x>0$, the following statements are equivalent:
\begin{itemize}
\item Every nontrivial zero of $\zeta$ has real part $1/2$.
\item Every zero of $s\mapsto\Zcal(s,x)$ inside $0<\Re s<1$ has real part $1/2$.
\end{itemize}
\textbf{Proof.} The first statement implies the second by Theorem 1. Conversely, by Theorem 2 every nontrivial zeta zero has a preimage with the same real part. If all these preimages lie on the critical line, then the original zeros lie there as well.

This equivalence is the essential consequence of the substitution. It imposes no additional restriction on the real parts and does not decide the Riemann Hypothesis.

\textbf{Why the region matters.} For $x\ge1$, one even has $f_a(x)\in(0,1)$ if and only if $a\in(0,1)$. Indeed, for positive integers $n$, $f_n(x)\ge1$, hence also for $a\ge1$; negative indices give nonpositive values. Thus, for $x\ge1$, all preimages of nontrivial zeta zeros lie in the input region $0<\Re s<1$.

For $0<x<1$, additional preimages outside this strip can occur. At $x=1/2$, for example, $f_2(x)=1/2$. Given a known zero $1/2+i\gamma$, Theorem 2 gives a $t$ with $h_{1/2}(t)=\gamma$, hence $\Zcal(2+it,1/2)=0$. The zeta argument still has real part $1/2$.

\section{The Height Map Can Produce Multiple Preimages}
On every interval $[n,n+1]$, the equation $h_x(t)=\gamma$ can be solved linearly. If
\[
 D_n(x)=f_{n+1}(x)-f_n(x)
\]
does not vanish, the solution is
\[
 t=n+\frac{\gamma-f_n(x)}{D_n(x)}.
\]
It lies in the interval exactly when $\gamma$ lies between the two endpoint values. If $D_n(x)=0$, either there is no solution or the entire interval consists of solutions.

The number of preimages therefore gives no new statement about the number of original zeta zeros.

\textbf{Concrete example $x=3$.} The support values at $n=7,8,9,10$ are $13,27,16,30$. Hence the first positive zeta-zero height has one positive preimage, whereas the second has three. The known heights are given in [4]:

\begin{center}
\begin{tabularx}{\linewidth}{@{}lX@{}}
\toprule
\textbf{Original height $\gamma$} & \textbf{Positive heights $t$ of the preimages at $x=3$}\\
\midrule
$14.1347251417\ldots$ & $7.0810517958\ldots$\\
$21.0220396388\ldots$ & $7.5730028313\ldots,\ 8.5434509419\ldots,\ 9.3587171171\ldots$\\
\bottomrule
\end{tabularx}
\end{center}

All associated indices here have real part $1/2$. The list is complete: for $t\ge0$, $h_3(t)\ge t$, because the positive support polynomials have nonnegative coefficients and $f_n(3)\ge f_n(2)=n$. Thus only finitely many intervals below the respective height need to be checked; their endpoint values were evaluated. The decimal values are numerical illustrations, not certified enclosures.

For $x=1$, every level $\gamma>1$ even has infinitely many positive preimages: $h_1(2^k)=1$, while $h_1(3^k)=2^k$ grows without bound. Arbitrarily far to the right there are therefore alternating values below and above the level. For $0<x<1$, the same holds for every $\gamma>0$, because $h_x(2^k)=x^k\to0$ and $h_x(3^k)\to\infty$. This shows how strongly the substitution changes the height ordering.

\section{For a Fixed Index, Zeros Are Pulled Back Polynomially}
\textbf{Theorem 4. Preimages and multiplicities in the variable $x$.}
Let $P(x)=f_s(x)$, for a fixed complex index $s$, be a nonconstant polynomial of degree $d$. Then $x\mapsto\zeta(P(x))$ is meromorphic. To every zeta zero $\rho$ there correspond exactly $d$ solutions of $P(x)=\rho$, counted with their polynomial multiplicities.

If $\zeta$ has a zero of order $m$ at $\rho$ and $P-\rho$ has a zero of order $k$ at $x_0$, then
\[
 \operatorname{ord}_{x_0}\zeta(P(x))=mk.
\]
\textbf{Proof.} Locally, $\zeta(w)=(w-\rho)^m g(w)$ with $g(\rho)\ne0$ and $P(x)-\rho=(x-x_0)^k q(x)$ with $q(x_0)\ne0$. Substitution gives the stated order. The number of polynomial preimages follows from the Fundamental Theorem of Algebra.

The only poles lie at $P(x)=1$. Their order is the multiplicity of the corresponding zero of $P-1$, because $\zeta$ has a simple pole at 1 [1]. If $P'(x_0)\ne0$, the residue is $1/P'(x_0)$. If $P$ is constant, this exceptional case must be treated separately; in particular, when $P\equiv1$ the composition is not a meromorphic function with isolated poles.

\textbf{Relation to the Pratt zeros.} A zero of $f_s$ itself gives $\zeta(f_s(x))=\zeta(0)=-1/2$ [4]. Thus the irreducibility and zero locations of the original polynomials concern equations different from $f_s(x)=\rho$. Those properties alone imply no critical-line statement for $\zeta$.

\textbf{Example 5. Three exact preimages of the first zeta zero.}
Let $\rho_1=1/2+i\gamma_1$ be the first positive zeta zero, with $\gamma_1=14.1347251417\ldots\in(14,15)$ [4]. The recursion gives
\[
 Q(x):=f_{14}(x)=x^3+x^2+x,\qquad f_{15}(x)=Q(x)+1.
\]
Hence, exactly and without having to determine the height algebraically,
\[
 f_{\rho_1}(x)=\frac12+i\bigl(Q(x)+\gamma_1-14\bigr),
\]
\[
 f_{\rho_1}(x)-\rho_1=i\bigl(Q(x)-14\bigr)
 =i(x-2)(x^2+3x+7).
\]
The three preimages of this one zeta zero are therefore
\[
 x=2,\qquad x=\frac{-3+i\sqrt{19}}2,\qquad x=\frac{-3-i\sqrt{19}}2.
\]
At all three points, $\Zcal(\rho_1,x)=0$. The polynomial preimages are simple, so the multiplicity of each composed zero is exactly the multiplicity of $\rho_1$. This factorization therefore does not prove simplicity of a zeta zero.

\section{Functional Equation and Analytic Limitation}
The completed zeta function is the entire function [3]
\[
 \xi(w)=\frac12 w(w-1)\pi^{-w/2}\Gamma(w/2)\zeta(w),\qquad
 \xi(w)=\xi(1-w).
\]
Its zeros are exactly the nontrivial zeta zeros; the apparent special points in the product representation are handled by continuation. For $0\le\Re s\le1$, even for complex $x$,
\[
 f_{1-s}(x)=1-f_s(x).
\]
Consequently, $\Xcal(s,x)=\xi(f_s(x))$ inherits in this strip the symmetry
\[
 \Xcal(1-s,x)=\Xcal(s,x).
\]
For real $x$, conjugation symmetry is added. Such a symmetry, however, also permits pairs of zeros away from the critical line; it forces no individual zero onto that line.

\textbf{The decisive analytic limitation.} For $x\ne2$, the map $s\mapsto\Zcal(s,x)$ is generally not holomorphic. On an open cell with $0<\sigma<1$ and $n<t<n+1$, outside the poles,
\[
 \frac{\partial\Zcal}{\partial\overline s}
 =\frac12\zeta'(f_s(x))\bigl(1-D_n(x)\bigr).
\]
\textbf{Proof.} On the cell, $\partial_\sigma f_s=1$ and $\partial_t f_s=iD_n$. The Wirtinger derivative $\partial_{\overline s}=(\partial_\sigma+i\partial_t)/2$ gives the formula by the chain rule. Already on $1<t<2$, $D_1=x-1$, so the additional factor is $(2-x)/2$. For $x\ne2$, the derivative does not vanish identically there, because $\zeta'$ does not vanish on an entire nonconstant interval of arguments.

At $x=2$, by contrast, we recover the original meromorphic function $\zeta(s)$. Analytic tools in the index variable, such as the argument principle, therefore do not transfer directly for other $x$. In the variable $x$, the meromorphic structure described in Theorem 4 remains intact.

\section{A Purely Geometric, Cubic Form of the Hypothesis}
A further exact reformulation needs only the polynomial $Q$ from Example 5. Choose the fixed index $s_0=1/2+14i$; this index need not itself be a zeta zero. Then
\[
 f_{s_0}(x)=\frac12+iQ(x),\qquad
 W(x)=\xi\!\left(\frac12+iQ(x)\right).
\]
The function $W$ is entire in $x$. Writing $x=u+iv$, we obtain
\[
 \Re\!\left(\frac12+iQ(u+iv)\right)
 =\frac12-v\bigl(3u^2-v^2+2u+1\bigr).
\]
The preimage of the critical line is therefore the following union of the real axis with two hyperbola branches:
\[
 C=\{u+iv:v=0\}\ \cup\ \{u+iv:v^2=3u^2+2u+1\}.
\]

\textbf{Theorem 6. Equivalence with a zero curve in the $x$-plane.}
The Riemann Hypothesis holds if and only if all zeros of $W$ lie on $C$.

\textbf{Proof.} If the hypothesis holds, then every zeta argument at a zero of $W$ has real part $1/2$, so by the formula above $x$ lies on $C$. Conversely, every nontrivial zeta zero $\rho$ has a preimage under the nonconstant polynomial $1/2+iQ$. If all zeros of $W$ lie on $C$, then this $\rho$ also has real part $1/2$.

For every real $y$, the equation $Q(x)=y$ has exactly one real and two conjugate nonreal solutions. Indeed, $Q'(u)=3u^2+2u+1>0$ on the real axis. The three solutions are simple: the critical values of $Q$ are $(-7\pm4i\sqrt2)/27$ and hence are not real. Under the Riemann Hypothesis, each nontrivial zeta zero is distributed over these three curve points, with its multiplicity unchanged at each point.

\begin{figure}[htbp]
\centering
\includegraphics[width=.72\textwidth]{pratt-zeta-curve-en.png}
\caption{Blue curve: the exact preimage of the critical line under $1/2+iQ(x)$. Markers: numerical preimages of the first conjugate pair of zeta zeros. The curves also contain points that are not zeros.}
\end{figure}

The curve is derived exactly. The six marked points are numerically computed preimages of the two known zeros $1/2\pm i\gamma_1$ under $1/2+iQ$. They illustrate the equivalence and do not replace a proof that all zeros lie on the curve.

\section{Result and the Concrete Remaining Task}
The substitution gives, first, a complete zero correspondence in the critical strip; second, explicit polynomial preimages with controlled multiplicities; and third, the cubic curve formulation of Theorem 6. A false Riemann Hypothesis would likewise appear in each of these equivalent pictures as a zero away from the corresponding line or curve.

For a proof, an independent property would now be needed that forces all zeros of
\[
 W(x)=\xi\!\left(\frac12+i(x^3+x^2+x)\right)
\]
onto $C$, or equivalently forces all zeros of $\Zcal$ in the input strip onto $\Re s=1/2$. Neither the interpolation nor integer multiplicativity nor the irreducibility of the prime-index polynomials has so far supplied such a property.

\clearpage
\section*{References}
\textbf{Verification.} The identities for $f_{14}$ and $f_{15}$, the factorization of $Q-14$, and the value $Q'(2)=17$ were checked using exact integer polynomial arithmetic. The decimal examples and curve points were computed numerically; supplementary Euler-Maclaurin evaluations reproduced residuals below $2\cdot10^{-14}$ at the first two known zeta zeros. This is a plausibility check, not a certified zero enclosure. The theorems rely on the stated proofs.

[1] NIST Digital Library of Mathematical Functions, \S25.2: definition and meromorphic continuation of the Riemann zeta function. \url{https://dlmf.nist.gov/25.2}

[2] NIST Digital Library of Mathematical Functions, \S25.10: trivial and nontrivial zeros, critical strip, and Riemann Hypothesis. \url{https://dlmf.nist.gov/25.10}

[3] NIST Digital Library of Mathematical Functions, \S25.4: functional equation and definition of the completed zeta function. \url{https://dlmf.nist.gov/25.4}

[4] Official mpmath documentation, Riemann zeta function: numerical values of the first zeros and $\zeta(0)$. \url{https://mpmath.readthedocs.io/en/latest/functions/zeta.html}

The definition of the Pratt family comes from the attached photograph and the user's supplementary specifications. This note studies the resulting composition; the proved zero correspondences do not require an irreducibility hypothesis.

\end{document}
